\documentclass[../../main.tex]{subfiles}
\begin{document}
\section{习题}\label{section:ch08-习题}

\begin{exercise}\label{exercise:ch08-01}
证明\cref{lemma:ch08-06}（向量场的延拓引理）。
\end{exercise}

\begin{exercise}\label{exercise:ch08-02}[Euler 齐次函数定理]
设 $c$ 是实数，$f\colon\mathbb{R}^n\setminus\{0\}\to\mathbb{R}$ 是光滑函数，关于次数 $c$ 正齐次，即对所有 $\lambda>0$ 和 $x\in\mathbb{R}^n\setminus\{0\}$ 有 $f(\lambda x)=\lambda^c f(x)$。证明 $Vf=cf$，其中 $V$ 是\cref{example:ch08-03}中定义的 Euler 向量场。（用于第 248 页。）
\end{exercise}

\begin{exercise}\label{exercise:ch08-03}
设 $M$ 是非空正维数的带边或不带边光滑流形。证明 $\mathfrak{X}(M)$ 是无限维的。
\end{exercise}

\begin{exercise}\label{exercise:ch08-04}
设 $M$ 是带边光滑流形。证明存在 $M$ 上的全局光滑向量场，其在 $\partial M$ 上的限制处处指向内侧，也存在一个在 $\partial M$ 上的限制处处指向外侧。（用于第 223、386 页。）
\end{exercise}

\begin{exercise}\label{exercise:ch08-05}
证明\cref{proposition:ch08-11}（局部标架的补全）。
\end{exercise}

\begin{exercise}\label{exercise:ch08-06}
设 $H$ 是四元数代数，$\Sigma\subseteq H$ 是单位四元数群（见\cref{exercise:ch07-22}和\cref{exercise:ch07-23}）。
\begin{enumerate}[(a)]
\item 证明若 $p\in H$ 是虚的，则 $qp$ 在每个 $q\in\Sigma$ 处与 $\Sigma$ 相切。（这里我们以通常方式将 $H$ 的每个切空间等同于 $H$ 本身。）

\item 在 $H$ 上定义向量场 $X_1,X_2,X_3$ 为
\begin{align*}
X_1|_q=qi,\qquad X_2|_q=qj,\qquad X_3|_q=qk.
\end{align*}
证明这些向量场限制为 $\Sigma$ 上的光滑左不变全局标架。

\item 在 $\mathbb{R}^4$ 与 $H$ 之间的同构 $(x^1,x^2,x^3,x^4)\leftrightarrow x^1+x^2\mathrm{i}+x^3j+x^4k$ 下，证明这些向量场有如下坐标表示：
\begin{align*}
X_1&=-x^2\frac{\partial}{\partial x^1}+x^1\frac{\partial}{\partial x^2}+x^4\frac{\partial}{\partial x^3}-x^3\frac{\partial}{\partial x^4},\\
X_2&=-x^3\frac{\partial}{\partial x^1}-x^4\frac{\partial}{\partial x^2}+x^1\frac{\partial}{\partial x^3}+x^2\frac{\partial}{\partial x^4},\\
X_3&=-x^4\frac{\partial}{\partial x^1}+x^3\frac{\partial}{\partial x^2}-x^2\frac{\partial}{\partial x^3}+x^1\frac{\partial}{\partial x^4}.
\end{align*}
\end{enumerate}
（用于第 179、562 页。）
\end{exercise}

\begin{exercise}\label{exercise:ch08-07}
八元数代数（也称为 Cayley 数）是八维实向量空间 $\mathbb{O}=H\oplus H$（其中 $H$ 是\cref{exercise:ch07-22}中定义的四元数空间），赋予如下双线性乘积：
\begin{align}
(p,q)(r,s)=\bigl(pr-s\overline{q},\,\overline{p}s+rq\bigr),\label{align:ch08-20}
\end{align}
对 $p,q,r,s\in H$。证明 $\mathbb{O}$ 是 $\mathbb{R}$ 上的非交换、非结合代数，并通过尽可能模仿\cref{exercise:ch08-06}证明 $S^7$ 上存在光滑全局标架。[提示：证明对所有 $P,Q\in\mathbb{O}$ 有 $(P\overline{Q})Q=P(\overline{Q}Q)$ 可能有帮助，其中 $(p,q)^*=(p^*,-q)$。关于八元数的更多信息，见 [Bae02]。]（用于第 179 页。）
\end{exercise}

\begin{exercise}\label{exercise:ch08-08}
十六元数代数是十六维实向量空间 $\mathbb{S}=\mathbb{O}\oplus\mathbb{O}$，乘积由\eqref{align:ch08-20}定义，但 $p,q,r,s$ 解释为 $\mathbb{O}$ 的元素。为什么十六元数乘法不能给出 $S^{15}$ 的全局标架？[注："sedenion"一名来自拉丁语 sedecim，意为十六。可除代数是具有乘法单位元且无零因子的代数（即 $ab=0$ 当且仅当 $a=0$ 或 $b=0$）。由 Bott、Milnor 和 Kervaire 关于球面可平行化性的工作 [MB58, Ker58] 可知，$\mathbb{R}$ 上的有限维可除代数维数必须为 $1$、$2$、$4$ 或 $8$。]
\end{exercise}

\begin{exercise}\label{exercise:ch08-09}
通过寻找反例证明，若将 $F$ 是微分同胚的假设替换为 $F$ 光滑且双射的较弱假设，则\cref{proposition:ch08-19}不成立。
\end{exercise}

\begin{exercise}\label{exercise:ch08-10}
设 $M$ 是 $\mathbb{R}^2$ 中 $x$ 和 $y$ 均为正的开子流形，$F\colon M\to M$ 为映射 $F(x,y)=(xy,y/x)$。证明 $F$ 是微分同胚，并计算 $F_*X$ 和 $F_*Y$，其中
\begin{align*}
X=x\frac{\partial}{\partial x}+y\frac{\partial}{\partial y},\qquad
Y=y\frac{\partial}{\partial x}.
\end{align*}
\end{exercise}

\begin{exercise}\label{exercise:ch08-11}
对平面上下列每个向量场，计算其在右半平面 $\{(x,y):x>0\}$ 上极坐标中的坐标表示。
\begin{enumerate}[(a)]
\item $X=x\dfrac{\partial}{\partial x}+y\dfrac{\partial}{\partial y}$。
\item $Y=x\dfrac{\partial}{\partial y}-y\dfrac{\partial}{\partial x}$。
\item $Z=(x^2+y^2)\dfrac{\partial}{\partial x}$。
\end{enumerate}
\end{exercise}

\begin{exercise}\label{exercise:ch08-12}
设 $F\colon\mathbb{R}^2\to\mathbb{RP}^2$ 为光滑映射 $F(x,y)=[x,y,1]$，$X\in\mathfrak{X}(\mathbb{R}^2)$ 由 $X=x\partial/\partial y-y\partial/\partial x$ 定义。证明存在 $Y\in\mathfrak{X}(\mathbb{RP}^2)$ 与 $X$ 是 $F$-相关的，并计算其在\cref{example:ch01-05}中定义的每个坐标卡下的坐标表示。
\end{exercise}

\begin{exercise}\label{exercise:ch08-13}
证明 $S^2$ 上存在恰好于一点处为零的光滑向量场。[提示：尝试使用球极投影；见\cref{exercise:ch01-07}。]
\end{exercise}

\begin{exercise}\label{exercise:ch08-14}
设 $M$ 是带边或不带边的光滑流形，$N$ 是光滑流形，$f\colon M\to N$ 是光滑映射。定义 $F\colon M\to M\times N$ 为 $F(x)=(x,f(x))$。证明对每个 $X\in\mathfrak{X}(M)$，存在 $M\times N$ 上的光滑向量场与 $X$ 是 $F$-相关的。
\end{exercise}

\begin{exercise}\label{exercise:ch08-15}[子流形上向量场的延拓引理]
假设 $M$ 是光滑流形，$S\subseteq M$ 是带边或不带边的嵌入子流形。给定 $X\in\mathfrak{X}(S)$，证明在 $M$ 中 $S$ 的邻域上存在光滑向量场 $Y$，使得 $X=Y|_S$。证明每个此类向量场可延拓到整个 $M$，当且仅当 $S$ 是真嵌入的。
\end{exercise}

\begin{exercise}\label{exercise:ch08-16}
对下列定义在 $\mathbb{R}^3$ 上的每对向量场 $X,Y$，计算 Lie 括号 $[X,Y]$。
\begin{enumerate}[(a)]
\item $X=y\dfrac{\partial}{\partial z}-2xy^2\dfrac{\partial}{\partial y}$；$Y=\dfrac{\partial}{\partial y}$。
\item $X=x\dfrac{\partial}{\partial y}-y\dfrac{\partial}{\partial x}$；$Y=y\dfrac{\partial}{\partial z}-z\dfrac{\partial}{\partial y}$。
\item $X=x\dfrac{\partial}{\partial y}-y\dfrac{\partial}{\partial x}$；$Y=x\dfrac{\partial}{\partial y}+y\dfrac{\partial}{\partial x}$。
\end{enumerate}
\end{exercise}

\begin{exercise}\label{exercise:ch08-17}
设 $M$ 和 $N$ 是光滑流形。给定向量场 $X\in\mathfrak{X}(M)$ 和 $Y\in\mathfrak{X}(N)$，可由
\begin{align*}
(X\oplus Y)_{(p,q)}=(X_p,Y_q),
\end{align*}
定义 $M\times N$ 上的向量场 $X\oplus Y$，其中我们将右端视为 $T_pM\oplus T_qN$ 的元素，它由\cref{proposition:ch03-14}自然等同于 $T_{(p,q)}(M\times N)$。证明若 $X$ 和 $Y$ 光滑，则 $X\oplus Y$ 光滑，且 $[X_1\oplus Y_1,X_2\oplus Y_2]=[X_1,X_2]\oplus[Y_1,Y_2]$。（用于第 346、527 页。）
\end{exercise}

\begin{exercise}\label{exercise:ch08-18}
假设 $F\colon M\to N$ 是光滑淹没，其中 $M$ 和 $N$ 是正维数光滑流形。给定 $X\in\mathfrak{X}(M)$ 和 $Y\in\mathfrak{X}(N)$，若 $X$ 和 $Y$ 是 $F$-相关的，则称 $X$ 是 $Y$ 的\textbf{提升}。向量场 $V\in\mathfrak{X}(M)$ 称为\textbf{竖直的}，若 $V$ 处处与 $F$ 的纤维相切（等价地，若 $V$ 与 $N$ 上的零向量场是 $F$-相关的）。
\begin{enumerate}[(a)]
\item 证明若 $\dim M=\dim N$，则 $N$ 上的每个光滑向量场都有唯一提升。
\item 证明若 $\dim M\neq\dim N$，则 $N$ 上的每个光滑向量场都有提升，但不唯一。
\item 再假设 $F$ 是满射。给定 $X\in\mathfrak{X}(M)$，证明 $X$ 是 $N$ 上某光滑向量场的提升，当且仅当每当 $F(p)=F(q)$ 时有 $\mathrm{d}F_p(X_p)=\mathrm{d}F_q(X_q)$。证明若如此，则 $X$ 是唯一光滑向量场的提升。
\item 再假设 $F$ 是满射且纤维连通。证明向量场 $X\in\mathfrak{X}(M)$ 是 $N$ 上某光滑向量场的提升，当且仅当每当 $V\in\mathfrak{X}(M)$ 竖直时 $[V,X]$ 竖直。
\end{enumerate}
（用于第 434 页。）
\end{exercise}

\begin{exercise}\label{exercise:ch08-19}
证明带叉积的 $\mathbb{R}^3$ 是 Lie 代数。
\end{exercise}

\begin{exercise}\label{exercise:ch08-20}
设 $A\subseteq\mathfrak{X}(\mathbb{R}^3)$ 是由 $\{X,Y,Z\}$ 张成的子空间，其中
\begin{align*}
X=y\frac{\partial}{\partial z}-z\frac{\partial}{\partial y},\quad
Y=z\frac{\partial}{\partial x}-x\frac{\partial}{\partial z},\quad
Z=x\frac{\partial}{\partial y}-y\frac{\partial}{\partial x}.
\end{align*}
证明 $A$ 是 $\mathfrak{X}(\mathbb{R}^3)$ 的 Lie 子代数，同构于带叉积的 $\mathbb{R}^3$。（用于第 538 页。）
\end{exercise}

\begin{exercise}\label{exercise:ch08-21}
证明在同构意义下，恰有一个一维 Lie 代数和两个二维 Lie 代数。证明这三个代数都同构于 $\mathfrak{gl}(2,\mathbb{R})$ 的 Lie 子代数。
\end{exercise}

\begin{exercise}\label{exercise:ch08-22}
设 $A$ 是 $\mathbb{R}$ 上的任意代数。$A$ 的\textbf{导子}是满足对所有 $x,y\in A$ 有 $D(xy)=(Dx)y+x(Dy)$ 的线性映射 $D\colon A\to A$。证明若 $D_1$ 和 $D_2$ 是 $A$ 的导子，则 $[D_1,D_2]=D_1\circ D_2-D_2\circ D_1$ 也是导子。证明 $A$ 的所有导子在该括号运算下构成 Lie 代数。
\end{exercise}

\begin{exercise}\label{exercise:ch08-23}
\begin{enumerate}[(a)]
\item 给定 Lie 代数 $\mathfrak{g}$ 和 $\mathfrak{h}$，证明直和 $\mathfrak{g}\oplus\mathfrak{h}$ 在由
\begin{align*}
\bigl[(X,Y),(X',Y')\bigr]=\bigl([X,X'],[Y,Y']\bigr)
\end{align*}
定义的括号下是 Lie 代数。

\item 假设 $G$ 和 $H$ 是 Lie 群。证明 $\operatorname{Lie}(G\times H)$ 同构于 $\operatorname{Lie}(G)\oplus\operatorname{Lie}(H)$。
\end{enumerate}
\end{exercise}

\begin{exercise}\label{exercise:ch08-24}
假设 $G$ 是 Lie 群，$\mathfrak{g}$ 是其 Lie 代数。向量场 $X\in\mathfrak{X}(G)$ 称为\textbf{右不变的}，若它在所有右平移下不变。
\begin{enumerate}[(a)]
\item 证明 $G$ 上右不变向量场的集合 $\overline{\mathfrak{g}}$ 是 $\mathfrak{X}(G)$ 的 Lie 子代数。

\item 设 $i\colon G\to G$ 表示取逆映射 $i(g)=g^{-1}$。证明前推 $i_*\colon\mathfrak{X}(G)\to\mathfrak{X}(G)$ 限制为从 $\mathfrak{g}$ 到 $\overline{\mathfrak{g}}$ 的 Lie 代数同构。
\end{enumerate}
\end{exercise}

\begin{exercise}\label{exercise:ch08-25}
证明若 $G$ 是交换 Lie 群，则 $\operatorname{Lie}(G)$ 是交换的。[提示：证明取逆映射 $i\colon G\to G$ 是群同态，并使用\cref{exercise:ch07-02}。]
\end{exercise}

\begin{exercise}\label{exercise:ch08-26}
假设 $F\colon G\to H$ 是 Lie 群同态。证明 $F_*\colon\operatorname{Lie}(G)\to\operatorname{Lie}(H)$ 的核是 $\operatorname{Ker}F$ 的 Lie 代数（在如\cref{theorem:ch08-46}中将子群的 Lie 代数等同为 Lie 子代数的意义下）。
\end{exercise}

\begin{exercise}\label{exercise:ch08-27}
设 $G$ 和 $H$ 是 Lie 群，$F\colon G\to H$ 是同时也是局部微分同胚的 Lie 群同态。证明诱导同态 $F_*\colon\operatorname{Lie}(G)\to\operatorname{Lie}(H)$ 是 Lie 代数同构。（用于第 531、557 页。）
\end{exercise}

\begin{exercise}\label{exercise:ch08-28}
将 $\det\colon\operatorname{GL}(n,\mathbb{R})\to\mathbb{R}^*$ 视为 Lie 群同态，证明其诱导 Lie 代数同态是 $\operatorname{tr}\colon\mathfrak{gl}(n,\mathbb{R})\to\mathbb{R}$。[提示：见\cref{exercise:ch07-04}。]
\end{exercise}

\begin{exercise}\label{exercise:ch08-29}
\cref{theorem:ch08-46}蕴涵 $\operatorname{GL}(n,\mathbb{R})$ 的任意 Lie 子群的 Lie 代数典范同构于 $\mathfrak{gl}(n,\mathbb{R})$ 的子代数，对 $\operatorname{GL}(n,\mathbb{C})$ 的 Lie 子群有类似陈述。在此同构下证明
\begin{align*}
\operatorname{Lie}\bigl(\operatorname{SL}(n,\mathbb{R})\bigr)&\cong\mathfrak{sl}(n,\mathbb{R}),\\
\operatorname{Lie}\bigl(\operatorname{SO}(n)\bigr)&\cong\mathfrak{o}(n),\\
\operatorname{Lie}\bigl(\operatorname{SL}(n,\mathbb{C})\bigr)&\cong\mathfrak{sl}(n,\mathbb{C}),\\
\operatorname{Lie}\bigl(U(n)\bigr)&\cong\mathfrak{u}(n),\\
\operatorname{Lie}\bigl(\operatorname{SU}(n)\bigr)&\cong\mathfrak{su}(n),
\end{align*}
其中
\begin{align*}
\mathfrak{sl}(n,\mathbb{R})&=\{A\in\mathfrak{gl}(n,\mathbb{R}):\operatorname{tr}A=0\},\\
\mathfrak{o}(n)&=\{A\in\mathfrak{gl}(n,\mathbb{R}):A^{T}+A=0\},\\
\mathfrak{sl}(n,\mathbb{C})&=\{A\in\mathfrak{gl}(n,\mathbb{C}):\operatorname{tr}A=0\},\\
\mathfrak{u}(n)&=\{A\in\mathfrak{gl}(n,\mathbb{C}):A^{*}+A=0\},\\
\mathfrak{su}(n)&=\mathfrak{u}(n)\cap\mathfrak{sl}(n,\mathbb{C}).
\end{align*}
\end{exercise}

\begin{exercise}\label{exercise:ch08-30}
通过给出显式同构，证明 $\mathfrak{su}(2)$ 和 $\mathfrak{o}(3)$ 是同构的 Lie 代数，且两者都同构于带叉积的 $\mathbb{R}^3$。
\end{exercise}

\begin{exercise}\label{exercise:ch08-31}
设 $\mathfrak{g}$ 是 Lie 代数。线性子空间 $\mathfrak{h}\subseteq\mathfrak{g}$ 称为 $\mathfrak{g}$ 中的\textbf{理想}，若每当 $X\in\mathfrak{h}$ 且 $Y\in\mathfrak{g}$ 时 $[X,Y]\in\mathfrak{h}$。
\begin{enumerate}[(a)]
\item 证明若 $\mathfrak{h}$ 是 $\mathfrak{g}$ 中的理想，则商空间 $\mathfrak{g}/\mathfrak{h}$ 有唯一的 Lie 代数结构，使得投影 $\pi\colon\mathfrak{g}\to\mathfrak{g}/\mathfrak{h}$ 是 Lie 代数同态。
\item 证明子空间 $\mathfrak{h}\subseteq\mathfrak{g}$ 是理想，当且仅当它是某个 Lie 代数同态的核。
\end{enumerate}
（用于第 533 页。）
\end{exercise}

\end{document}